\chapter{The pre-analysis plan and its log}\label{app:plan}

\section{The plan}
Before any outcome (basis, returns, drift, delay) was related to any treatment (ban status, utilisation, MWPL changes, STT
dates, F\&O entry or exit, the tax change), I wrote a pre-analysis plan fixing the samples, the outcome and treatment
definitions, the specifications, the predicted signs, the falsifiers, the inference and the primary and secondary
hypotheses. It was frozen with a time-stamped tag in the project's version control on 5 November 2026, after the guide's
go-ahead and a full audit of the data. Until then, only treatment-side counts and outcome distributions on their own had
been examined.

\begin{itemize}
\item \emph{Primary family:} H7 (one-sided correction), H8 (trading), H3 (costs and band edges), H6 (F\&O entry and exit),
corrected jointly by the Holm method at 5\%.
\item \emph{Secondary tests:} the level effect, H9, the regression discontinuity at 95\% with a density test, hysteresis,
H2, H1, H4 and H5, reported without adjustment and labelled exploratory.
\item \emph{Inference:} standard errors clustered by stock and date; bootstrap where clusters are few; winsorisation at 1\%
and 99\% within each year, with robustness without.
\end{itemize}

\section{The log}
Where the plan left a detail open (for example the exact sample filters, the composition split for the level effect, the
population for the density test, or the outcome window for drift), I wrote the detail into a dated log, committed it,
and only then ran the test. Analyses added after seeing a result are marked as exploratory in the text and listed in the
log; the pre-specified result is always reported first. The main entries were:
\begin{itemize}
\item the analysis sample, the H7 regimes and test statistic, the H8 comparison zone and identity reading, the H3 windows
and placebo, and the H6 design (calendar-year cohorts, controls, covariates, outcome definitions, stacked design for drift,
a Rs 5 price filter for event outcomes), all before the primary tests;
\item before each secondary test, its design: event and control definitions and a median split for the level effect (only
8\% of entries have falling open interest before entry); a density test on stocks below the trigger the day before
(because the density over all non-ban days thins above the trigger by construction); pre-determined placebo outcomes for
the discontinuity; five basis lags for hysteresis; stock-specific tail measures and the reform intensity for H9; the
interaction design for H2; a principal-component friction index for H1; lagged flows with tax and STT dummies for H4;
constraint groups for H5;
\item before the band estimation and the constraint index, their specifications;
\item exploratory additions: the enforcement-period split for H8, diagnostics and a noise correction for H7, the placebo
calibration of H3, an alternative baseline and a build-up classification for the level effect, bandwidth checks for the
density test, period-specific demeaning and a noise check for the band.
\end{itemize}
